%! Change in Tandem — Unit 1, Lesson 1.1 — Exit Ticket
\documentclass[10pt]{article}
\usepackage{precalculus-article}
\usepackage{precalculus-boxes}
\newcommand{\TallMath}[1]{$\displaystyle #1\rule[-1.4em]{0pt}{3.2em}$}

\begin{document}

\pageheader{Unit 1, Lesson 1.1}{Exit Ticket}
\namedateperiod

\vspace{0.14in}

\begin{enumerate}[itemsep=12pt]

\item A function $f$ is described as follows: as $x$ increases from 1 to 3, $f(x)$ increases
from 2 to 8. As $x$ increases from 3 to 6, $f(x)$ decreases from 8 to 1.

\begin{enumerate}[label=(\alph*), itemsep=6pt]
  \item State the domain and range of $f$ based on this description.

        Domain: \blank{4cm} \qquad Range: \blank{4cm}

  \item On what interval is $f$ \textbf{increasing}? \blank{3cm} \quad
        On what interval is $f$ \textbf{decreasing}? \blank{3cm}

  \item At what input value does $f$ reach its maximum output? \blank{3cm}
        What is that maximum output? \blank{3cm}
\end{enumerate}

\item The graph of $g$ is \textbf{concave up} on $[0, 4]$ and \textbf{concave down} on $[4, 8]$.

What does this tell you about how the \textit{rate of change} of $g$ is behaving on each interval?
Write one sentence for each interval.

\vspace{0.15in}
On $[0, 4]$: \writeline

\vspace{0.12in}
On $[4, 8]$: \writeline

\end{enumerate}

\end{document}
